\begin{claim}{C.92}{Proved}{Arbitrary pointed descents and the exact original lift topology}
Retain all the actual data of C.81: the original closed disk
$\overline D$, its interior $O$, the compatible closed-disk metric $d$,
the closed discrete free proper group $Q$ of operators on $O$,
the boundary-fixed centralizer $\mathscr C_Q$ with a transitive actual
subgroup, and the specified quotient covering $\pi:O\to N=O/Q$.
The quotient $N$ is allowed to be noncompact throughout this note.
Fix $x\in O$, put $p=\pi x$, and set $\mathscr F=\mathscr C_Q(x)$.
Write $\mathscr I_p$ for the identity path component of
$\operatorname{Homeo}(N\mathbin{\mathrm{rel}}\{p\})$ in the compact-open
topology of maps and inverses. Thus a path in this definition has
jointly continuous maps and inverses on $[0,1]\times N$.

Every member of the entire original $\mathscr F$ descends into
$\mathscr I_p$. Below we prove this assertion without assuming
component membership upstairs. We also identify its exact image
inside $\mathscr I_p$, with inverse maps and both relevant topologies.
The original compact-quotient contraction C.91 remains unchanged.


More precisely, let $\lambda(b)$ be the unique lift of $b\pi$
fixing $x$ for $b\in\mathscr I_p$. The image of $D$ is exactly
the subgroup $\mathscr B_p$ defined by
$\delta_b(\pi z)=d_Q(\lambda(b)z,z)\in C_0(N)$.
Descent and extension of this lift are inverse group maps and
are isometric for the original sum metric and the lifted sum
metric in (C.92.15). On a compact parameter space their exact
continuity criterion is the uniform tail (C.92.17).
The explicit noncompact example (C.92.18)--(C.92.21) gives a
sequence in this exact image converging in the quotient
compact-open topology whose pointed lifts stay more than
$2/5$ from the identity in the original metric.

\end{claim}

\begin{proof}
\emph{Literal descent and the original orbit labels.}
For $h\in\mathscr F$, define
\[
 D(h)(\pi z)=\pi(hz),\qquad
 D(h)^{-1}(\pi z)=\pi(h^{-1}z).
                                                               \tag{C.92.1}
\]
Commutation with every actual $q$ makes these formulas well-defined.
The quotient property makes them continuous, and substitution proves
they are inverse maps. They fix $p$. Composition gives
$D(hk)=D(h)D(k)$. Two such maps with equal descents are lifts of the
same map on the connected $O$, agree at $x$, and hence agree everywhere;
their boundary values are both the identity. Thus $D$ is injective.

For a loop $\gamma$ at $p$, let its lift starting at $x$ end at $qx$.
The image under $h$ of that lift starts at $x$ and ends at
$h(qx)=qh(x)=qx$. It lifts $D(h)\gamma$.
Since $O$ is simply connected, the two lifted paths with these
endpoints are homotopic rel endpoints. Their projections show
$D(h)_*[\gamma]=[\gamma]$ in the specified $\pi_1(N,p)$.
In particular every original orbit label is retained.

\emph{A based homotopy on the locally finite lifted triangulation.}
Here is a direct homotopy, including its maps, rather than an
inference of pointed isotopy from the preceding loop calculation.
Choose a disk coordinate $\theta:\overline D\to\{|u|\leq1\}\subset
\mathbb C$, put $a=\theta(x)$, and use the mutually inverse maps
\[
 \begin{aligned}
 m_a(u)&=(u-a)/(1-\overline a u),
 &m_a^{-1}(v)&=(v+a)/(1+\overline a v),\\
 \psi(z)&=\frac{m_a(\theta z)}{1-|m_a(\theta z)|},
 &\psi^{-1}(w)&=\theta^{-1}m_a^{-1}\left(\frac{w}{1+|w|}\right).
 \end{aligned}                                                  \tag{C.92.2}
\]
They give $\psi:O\to\mathbb R^2$ with $\psi(x)=0$.
The original covering and every operator $q$ remain the ones in
(C.92.1).

Use C.82's locally finite triangulation of the quotient surface,
refined locally so that $p$ is a vertex, and lift it by $\pi$.
Every lifted closed simplex maps homeomorphically to its quotient
simplex. The lifted family is locally finite, by the local covering
charts. Any deck operator preserving a lifted simplex fixes it
pointwise, by injectivity of the projection on that simplex, and is
therefore the identity by freeness. The lifts of a given quotient
simplex form one $Q$-orbit: match one point in two lifts by the unique
operator in its fiber and then use uniqueness of the lift on the
connected simplex. Choose one lift of each simplex and choose $x$
as the lift of the vertex $p$.

For a chosen vertex $v$ define
$W_t(v)=\psi^{-1}((1-t)\psi(hv)+t\psi(v))$, and define
$W_t(qv)=qW_t(v)$ on all translates. The path at $x$ is constant.
For a chosen edge, its endpoint paths and its prescribed maps
$h$ and the identity at times zero and one give a map $f$ on the
boundary of its square prism. For a chosen triangle the already
defined edge homotopies and the two end faces likewise give a map
on the boundary of its triangular prism.

For completeness, use the following explicit extension on either
convex coordinate prism $P$. Choose an interior point $c$. Every
point other than $c$ is uniquely $c+r(u-c)$ with $u\in\partial P$
and $0<r\leq1$. Define
\[
 \operatorname{Cone}_f(c+r(u-c))
       =\psi^{-1}\bigl(r\psi(f(u))\bigr),
 \qquad \operatorname{Cone}_f(c)=x.                            \tag{C.92.3}
\]
These radial coordinates are continuous away from $c$: on each
ray the boundary parameter is the first intersection with one
of the finitely many facet hyperplanes of the convex prism.
Equivalently, for a unit vector $v$ and inequalities
$a_i\cdot y\leq b_i$ defining $P$, that parameter is
$\min_{a_i\cdot v>0}(b_i-a_i\cdot c)/(a_i\cdot v)$.
This is positive and continuous on the unit sphere.
Continuity of (C.92.3) at $c$ follows from boundedness of
$\psi(f(\partial P))$. The value on $\partial P$ is exactly $f$.

Apply (C.92.3) first to each chosen edge prism and then to each
chosen triangle prism, using the coordinates of the triangulation,
and transport each extension by the literal operator $q$.
The boundary prescriptions agree on common faces. Local finiteness
gives continuity by closed-cell pasting on $O\times[0,1]$.
Thus we have actual maps
\[
 \begin{gathered}
 W:O\times[0,1]\to O,\qquad W_0=h|_O,\quad W_1=\id_O,\\
 W_t(qz)=qW_t(z),\qquad W_t(qx)=qx,\\
 B_t(\pi z)=\pi W_t(z),\qquad
 B_0=D(h),\quad B_1=\id_N,\quad B_t(p)=p .
 \end{gathered}                                                \tag{C.92.4}
\]
The last family is continuous by the local sections of $\pi$.
This is a based homotopy for every original $h\in\mathscr F$.
The argument uses local finiteness and does not require the
triangulation to be finite.

\emph{The orientation maps, including a nonorientable quotient.}
We verify every orientation qualification in the source theorem.
The connecting map
\[
 \partial:H_2(\overline D,\partial\overline D;\mathbb Z)
             \longrightarrow H_1(\partial\overline D;\mathbb Z)
                                                               \tag{C.92.5}
\]
is an isomorphism by the exact sequence of the contractible disk.
Naturality and $h|_{\partial\overline D}=\id$ give
$\partial h_*=\partial$, so $h_*$ is the identity on the left group.

Let $j_x$ be the map from
$H_2(\overline D,\partial\overline D)$ to
$H_2(\overline D,\overline D\setminus\{x\})$ induced by inclusion
of pairs. It is an isomorphism. Indeed, in the coordinate
$\kappa=m_a\theta$, the formula
\[
 R_s(z)=\kappa^{-1}
       \left(\left(1-s+\frac{s}{|\kappa(z)|}\right)\kappa(z)\right),
 \quad z\ne x,\quad 0\leq s\leq1,                             \tag{C.92.6}
\]
deformation retracts $\overline D\setminus\{x\}$ onto the boundary
while fixing the boundary. The natural diagram of the two exact
sequences then proves the assertion about $j_x$.
Let
$e_x:H_2(O,O\setminus\{x\})\to
H_2(\overline D,\overline D\setminus\{x\})$
be the excision isomorphism. Naturality for $h$, which fixes $x$,
gives, with $\eta_x=e_x^{-1}j_x$,
\[
 h_{x*}=\eta_x\,h_*\,\eta_x^{-1}
       =\id\quad\hbox{on }H_2(O,O\setminus\{x\};\mathbb Z).
                                                               \tag{C.92.7}
\]

Choose one specified evenly covered neighborhood $V_0$ of $p$
and its sheet $U_0$ containing $x$. Write $i_U,i_V$ for the
excision maps from the local homology groups of $U_0,V_0$
to those of $O,N$, respectively. The precise orientation
comparison at $x$ and $p$ is the isomorphism
\[
 \omega_x=i_V\circ(\pi|_{U_0})_*\circ i_U^{-1}:
 H_2(O,O\setminus\{x\})
        \longrightarrow H_2(N,N\setminus\{p\}).
                                                               \tag{C.92.8}
\]
Choose a smaller open neighborhood
$U\subset U_0\cap h^{-1}(U_0)$ of $x$.
On $\pi U$ the literal local formula is
\[
 D(h)|_{\pi U}
     =(\pi|_{U_0})\circ h|_U\circ(\pi|_U)^{-1}.
                                                               \tag{C.92.9}
\]
Excision and naturality applied to this formula prove
$D(h)_{p*}=\omega_xh_{x*}\omega_x^{-1}=\id$.
Thus the local orientation at the fixed point is preserved.
The use of a sheet in (C.92.8) is essential to the definition:
the global map $\pi$ sends every other point $qx$ in this fiber
to $p$, so the required local comparison is the displayed
excision composite.

On a connected orientable $N$, the sign of a homeomorphism on
the continuously chosen local fundamental classes is locally
constant and hence constant. Its sign at $p$ is $+1$ by
(C.92.7)--(C.92.9), so $D(h)$ preserves orientation globally.
This verifies the plane and open-cylinder cases.
On an open M\"obius band the same local-homology identity
verifies preservation of orientation near the basepoint.
Each point of $N$ has an open-disk covering chart, so
$\partial N=\varnothing$; the closed-cylinder and closed
M\"obius-band cases in the source statement therefore do
not occur as this $N$. No orientation of a nonorientable
$N$, and no orientation assumption on the original $q$,
has entered the local calculation.

\emph{The exact Epstein result and the topology of the resulting path.}
Use \cite[Theorem~6.4(b), p.~101; proof, pp.~101--103]{C_Epstein1966}.
Its based statement is the following: a self-homeomorphism of a
surface whose boundary components are compact, based homotopic
to the identity, is isotopic to the identity fixing that basepoint,
with orientation preservation required for the plane and the open
or closed cylinder, and local orientation preservation at the
basepoint required for the open or closed M\"obius band.
Equations (C.92.4) and (C.92.7)--(C.92.9) verify these hypotheses;
the boundary-component requirement holds because $\partial N$
is empty.

The exceptional M\"obius cases are proved in Lemma~6.5 on p.~102.
For the general case the proof on p.~103 concatenates isotopies
fixing successively larger compact submanifolds and explicitly
constructs a level-preserving homeomorphism of $N\times[0,1]$.
The topological version is included in the statement. Appendix
Theorem~A4 on p.~107 supplies its basepoint-fixed passage to a
PL homeomorphism; the corresponding topological proof paragraph
on p.~101 records the relative approximation step.
The hypotheses of Theorem~6.4(b) are met by the displayed
based homotopy itself; a properness assertion for (C.92.4)
on a noncompact $N$ is not used.

There is a small topology point in using this source:
Definition~1.4 on p.~84 defines an isotopy of embeddings.
The proof used here supplies homeomorphisms throughout.
For the general case this is explicit on p.~103.
The plane and open-cylinder steps in Theorems~5.4 and 5.7,
pp.~96--97, use Theorem~5.1, p.~95: its level maps for
$t<1$ are conjugates of the displayed homeomorphism, and
at $t=1$ are the identity. The compact disk step is the
explicit level-preserving homeomorphism of Theorem~5.2,
pp.~95--96. The open-M\"obius proof in Lemma~6.5 uses a
compact M\"obius core and the same half-cylinder construction
on its complement. In the cylinder step our based homotopy
also retains the generator of its fundamental group, as
proved after (C.92.1). After the based core-circle adjustment
in the source proof, the map fixes the core pointwise;
local orientation preservation makes it preserve its two
sides, so the two half-cylinder constructions apply to
the actual complementary components.
The topological approximation from A4 is an ambient
isotopy and likewise consists of homeomorphisms.

Write the resulting level-preserving homeomorphism as
$(y,t)\mapsto(b_t(y),t)$, where $b_0=D(h)$, $b_1=\id_N$,
and $b_t(p)=p$. Its inverse is exactly
$(y,t)\mapsto(b_t^{-1}(y),t)$ and is continuous.
This inverse continuity can also be checked without a
compactness assumption on $N$: extend the homeomorphism
family constantly for $t<0$ and $t>1$. Its level map is
a continuous bijection of the boundaryless $3$-manifold
$N\times\mathbb R$. Invariance of domain makes it open,
so its inverse is continuous; restrict back to $[0,1]$.
The compact-open exponential law on the locally compact
surface $N$ therefore gives a continuous path of maps and
inverses. Consequently $D(h)\in\mathscr I_p$ for every
$h\in\mathscr F$. This also places each such descent in the
pointed identity component; no equality of all components
or all possible descents has been presumed.

\emph{The exact pointed lift and its compact-open continuity.}
For a pointed homeomorphism $b$ of $N$, define the actual path
and its specified lift by
\[
 \beta_{b,z}(t)=b\!\left(\pi\psi^{-1}(t\psi(z))\right),
 \qquad
 \lambda(b)(z)=\widetilde\beta_{b,z}(1),\quad
 \widetilde\beta_{b,z}(0)=x.
                                                               \tag{C.92.10}
\]
The covering homotopy property applied with parameter $(b,z)$
gives joint continuity of the lifted paths. Indeed evaluation
of $b$ is continuous for the compact-open topology because
$N$ is locally compact, and the displayed formula is jointly
continuous before lifting. It gives
$\pi\lambda(b)=b\pi$ and $\lambda(b)(x)=x$.
Uniqueness of lifts on the connected $O$ makes it the unique
lift with that value. The same uniqueness proves the exact
identities
\[
 \lambda(bc)=\lambda(b)\lambda(c),\qquad
 \lambda(b^{-1})=\lambda(b)^{-1},\qquad
 \lambda(\id_N)=\id_O.
                                                               \tag{C.92.11}
\]
For the inverse identity, the two composites lift $\pi$
and fix $x$, and so are the identity. Thus these are actual
homeomorphisms. Applying the joint argument in (C.92.10) to
$b^{-1}$ also proves joint inverse continuity.
The exponential law on $O$ proves that $\lambda$ is continuous
from the pointed compact-open map-and-inverse topology to
the compact-open map-and-inverse topology on $O$.

For $b\in\mathscr I_p$, take a pointed path $b_t$ from the
identity to $b$. The continuous path $\lambda(b_t)(qx)$
lies in the discrete fiber $Qx$ and starts at $qx$.
It is therefore constantly $qx$. Now
$\lambda(b)q$ and $q\lambda(b)$ lift the same map $b\pi$
and agree at $x$, so uniqueness gives
\[
                  \lambda(b)q=q\lambda(b)
                         \quad\text{for every original }q.
                                                               \tag{C.92.12}
\]
This proves the individual operator equations without choosing
a deck representative modulo any subgroup.

Descent is continuous in the other direction as well.
The map $(h,z)\mapsto\pi(hz)$ on $\mathscr F\times O$
is continuous and is constant on the fibers of
$\id_{\mathscr F}\times\pi$. This product map is an open
quotient map, since $\pi$ is open. It therefore induces
the continuous joint map $(h,y)\mapsto D(h)y$.
Repeat the argument with $h^{-1}$ and use the exponential
law on $N$. This proves continuity of
$D:\mathscr F\to\mathscr I_p$ in the quotient compact-open
map-and-inverse topology.

\emph{The exact original image and its isometric topology.}
Retain $d_Q(z,w)=\sup_{q\in Q}d(qz,qw)$ from C.81 and put
\[
 \begin{split}
 \delta_b(\pi z)&=d_Q(\lambda(b)z,z),\\
 \mathscr B_p&=\{b\in\mathscr I_p:\delta_b\in C_0(N)\},\\
 L(b)|_O&=\lambda(b),\qquad
 L(b)|_{\partial\overline D}=\id\quad(b\in\mathscr B_p).
 \end{split}                                                   \tag{C.92.13}
\]
The displacement is well-defined and continuous by (C.92.12)
and $Q$-invariance of $d_Q$. C.88 gives its original boundary
extension exactly when this displacement belongs to $C_0(N)$.
For $h\in\mathscr F$, uniqueness gives
$\lambda(D(h))=h|_O$, so C.88 puts $D(h)$ in $\mathscr B_p$.
Conversely C.88 extends every $\lambda(b)$ in (C.92.13) to
an element of $\mathscr F$, and its descent is $b$. Thus
\[
 D:\mathscr F\longrightarrow\mathscr B_p,\qquad
 L:\mathscr B_p\longrightarrow\mathscr F,\qquad
 LD=\id_{\mathscr F},\quad DL=\id_{\mathscr B_p}
                                                               \tag{C.92.14}
\]
are the exact inverse group maps.
For example the subgroup assertion also follows directly from
$\delta_{bc}\leq\delta_b\circ c+\delta_c$ and
$\delta_{b^{-1}}=\delta_b\circ b^{-1}$: a homeomorphism and its
inverse carry compact sets to compact sets, so these pullbacks
preserve $C_0(N)$.

Define a metric on this exact image by the actual lifts:
\[
 d^\uparrow(b,c)=
 \sup_{z\in O}d_Q(\lambda(b)z,\lambda(c)z)
 +\sup_{z\in O}d_Q(\lambda(b)^{-1}z,\lambda(c)^{-1}z).
                                                               \tag{C.92.15}
\]
It is finite since $d$ is bounded on the compact original disk.
It is a metric by injectivity of the lift. Equivariance,
reindexing, and density of the interior give
\[
 \begin{split}
 \sup_{z\in O}d_Q(\lambda(b)z,\lambda(c)z)
 &=\sup_{z\in O,\ q\in Q}d(\lambda(b)qz,\lambda(c)qz)\\
 &=\sup_{w\in\overline D}d(L(b)w,L(c)w).
 \end{split}                                                   \tag{C.92.16}
\]
The same identity for the inverse maps proves that (C.92.14)
are isometric topological group isomorphisms for the original
closed-disk sum metric and $d^\uparrow$.
In particular the identity from this lifted metric on
$\mathscr B_p$ to its quotient compact-open topology is
continuous, by the descent argument above.

For a compact metric parameter space $P$ and a continuous
map $u\mapsto b_u$ into $\mathscr B_p$ with its quotient
compact-open topology, (C.92.10) gives joint interior continuity
of the actual pointed lifts and their inverses. C.88 now gives
the exact criterion for $u\mapsto L(b_u)$ to be continuous
in the original topology:
\[
 \forall\epsilon>0\ \exists K\subset N\text{ compact}\quad
 \sup_{\substack{u\in P\\y\notin K}}\delta_{b_u}(y)<\epsilon .
                                                               \tag{C.92.17}
\]
Its inverse family then satisfies the same uniform tail
condition. Thus the compact-family requirement is retained
on the unique lifts in (C.92.10), with every deck label fixed.

When $N$ is compact, taking $K=N$ makes (C.92.17) automatic,
and $C_0(N)=C(N)$ in the definition used here.
Then $\mathscr B_p=\mathscr I_p$.
For a convergent sequence $b_n\to b$ in the quotient topology,
apply (C.92.17) to $\{0\}\cup\{1/n:n\geq1\}$ and the lifts from
(C.92.10). This gives $L(b_n)\to L(b)$ in the original metric.
Both topologies are metrizable, so this proves continuity
of the inverse comparison in the compact case.
C.91's pointed-component contraction makes that entire
component path connected, giving
$\mathscr I_p=\operatorname{Homeo}(N\mathbin{\mathrm{rel}}\{p\})^0$
there. These are precisely C.91's original maps and compact
hypothesis.

\emph{An exact noncompact sequence, including its deck labels.}
The following actual C.81 example records why the last compact
argument cannot be replaced just by quotient compact-open
convergence. In this example $\overline D$ is the closed unit disk with its
Euclidean metric $d$. Use the original radial disk coordinate
\[
 j:\mathbb R^2\to\{|w|<1\},\quad j(z)=\frac{z}{1+|z|},
 \qquad j^{-1}(w)=\frac{w}{1-|w|},
                                                               \tag{C.92.18}
\]
and take $Q=\langle jTj^{-1}\rangle$ with
$T(s,t)=(s+1,t)$. The quotient coordinate is
$\pi(j(s,t))=([s],t)$ on $N=(\mathbb R/\mathbb Z)\times\mathbb R$.
The actual extended translations $j(z\mapsto z+a)j^{-1}$
are transitive, fix the original circle, and commute with
every displayed operator. The action is closed, discrete,
free and proper by the explicit plane coordinate, so this
is an instance of the retained hypotheses.

For an integer $R\geq2$ put
$\beta(v)=\max(0,1-4|v|)$ and
\[
 F_{R,t}(s,y)=(s+tR\beta(y-R),y),\qquad
 F_{R,t}^{-1}(s,y)=(s-tR\beta(y-R),y),
 \quad 0\leq t\leq1 .
                                                               \tag{C.92.19}
\]
They commute with every $T^n$, fix $(0,0)$, and have quotient
support in
$(\mathbb R/\mathbb Z)\times[R-\tfrac14,R+\tfrac14]$.
Both displacements in the plane are bounded by $R$.
For $|w|\leq R$ and $|z|>R$, direct subtraction gives
\[
       |j(z+w)-j(z)|\leq\frac{2R}{1+|z|-R}.
                                                               \tag{C.92.20}
\]
Thus $jF_{R,t}j^{-1}$ and its inverse extend jointly over
$[0,1]$ by the identity on the original circle.
Let $b_R$ be the descent of this endpoint.
Then $b_R\in\mathscr B_p$, with $x=j(0,0)$.
Its unique pointed lift is exactly $jF_{R,1}j^{-1}$.
Both $b_R$ and $b_R^{-1}$ are eventually the identity
on any specified compact subset of $N$, so $b_R\to1$
in the quotient compact-open map-and-inverse topology.

At $y_R=([0],R)$ the quotient value is $b_R(y_R)=y_R$,
because the horizontal translation there is the integer $R$.
The actual pointed lift, however, sends
$j(0,R)$ to $j(R,R)$, and therefore
\[
 \begin{gathered}
 \rho(b_R(y_R),y_R)=0,\qquad
 \delta_{b_R}(y_R)\geq
       |j(R,R)-j(0,R)|\geq\frac{R}{1+\sqrt2R}>\frac25,\\
 d^\uparrow(b_R,1)>\frac25 .
 \end{gathered}                                                \tag{C.92.21}
\]
For the strict scalar bound, put $R=1+s$, $s\geq0$.
Then $5R-2>0$ and
$(5R-2)^2-8R^2=1+14s+17s^2>0$.
The point-fixing lift in (C.92.10) cannot be replaced by a
different deck translate: a nonidentity translate sends $x$
to a different point of its fiber.
This is a sequence on the exact image $\mathscr B_p$ for
which the inverse compact-open comparison fails to be continuous.
It also fails the compact-family tail (C.92.17) on the
convergent-sequence parameter space, since the $y_R$ escape
all compact sets.

The proved comparison is consequently the original-image
description (C.92.13)--(C.92.17), together with pointed isotopy
of every actual descent. A deformation of a noncompact
quotient family still requires (C.92.17) for the actual
lift and its inverse throughout the deformation.
No such deformation is constructed in this note.
C.91 continues to assert the whole-group contraction under
its existing compact-quotient hypothesis.

\end{proof}
